\section{Definición del proyecto}
{\consigna
Esta sección debe definir el proyecto. A diferencia de la introducción, que es más ``blanda'', a lo largo de esta sección y subsecciones, se debe ser claro y específico, tratando de minimizar posibles ambigüedades.
Si vamos a hacer un sistema compuesto por un microcontrolador, sensores para medir ciertos indicadores, registrarlos y decidir tomar alguna acción en base a esos datos con cierta finalidad, digamos eso mismo: ``un sistema compuesto por un microcontrolador, sensores para medir ciertos indicadores, registrarlos y decidir tomar alguna acción en base a esos datos con'' la finalidad que sea.\\


En esta parte introductoria de la sección, se puede describir el objetivo de negocio que se persigue y cómo se espera cumplirlo, de forma de facilitar el entendimiento de las siguientes subsecciones.
}

El proyecto consiste en el desarrollo de un sistema portátil e inteligente compuesto por dos dispositivos complementarios —un dispositivo no invasivo ubicado en la parte posterior del brazo, formado por un parche sensor descartable y un módulo electrónico reutilizable de lectura y comunicación que se acopla a él, y una pulsera de muñeca reutilizable— junto con una capa de procesamiento e inteligencia artificial, destinado a personas con diabetes tipo 2 y pre-diabetes.

El parche sensor descartable permite medir glucosa, lactato y cortisol en sudor mediante el módulo electrónico reutilizable, mientras que la pulsera capta señales fisiológicas de contexto (frecuencia cardíaca, actividad física y temperatura de piel). El sistema integra ambas fuentes de datos para no solo reportar el nivel de glucosa, sino explicar su causa probable (alimentación, estrés, actividad física o una posible complicación renal) y predecir su evolución en dos modos complementarios: corto plazo (30 a 60 minutos), orientado a priorizar la exactitud de la predicción, y horizonte extendido (1 a 8 horas), orientado a estimar tendencias para apoyar la planificación del seguimiento, entregando al usuario una alerta temprana junto con una recomendación de acción concreta.

El dispositivo busca integrar en una única solución información que hoy se obtiene mediante distintos controles y análisis dispersos en el tiempo, complementando —no reemplazando— a los medidores de glucosa continuos (CGM) existentes en el mercado.


% la V 2.6 se reodena la sección misión y vision antes de Objetivo
\subsection{Misión y visión}
{\consigna
Esta subsección debe definir la Misión y la Visión del proyecto.\\
La **misión** es la razón de ser del proyecto basada en su propósito principal, mientras que la **visión** es la imagen de cómo se espera que el proyecto pueda evolucionar a futuro. La misión se enfoca en el presente y en lo que el resultado del proyecto hace para cumplir su propósito, mientras que la visión se enfoca en el futuro y en los posibles objetivos que el proyecto busca alcanzar.\\
Existe mucho material disponible para ayudar al alumno a definir estas dos características de su proyecto.
}

Desarrollar una solución portátil que integre mediciones no invasivas de biomarcadores en sudor y señales fisiológicas para aportar información contextual e identificar factores probablemente asociados a las variaciones observadas, sin atribuir causalidad clínica, con el propósito de facilitar el seguimiento de la salud metabólica.

A futuro, se busca consolidar una forma de seguimiento de la salud metabólica que combine mediciones, estimaciones e información contextual con evidencia suficiente para orientar su interpretación. En versiones posteriores podría evaluarse la incorporación de funciones relacionadas con complicaciones asociadas a la diabetes, como el daño renal, únicamente si existe evidencia que las sustente y se realiza una validación específica. La eventual contribución a la autonomía y seguridad de los pacientes o a un uso más eficiente de los recursos sanitarios deberá demostrarse mediante evaluaciones apropiadas.


\subsection{Objetivo}
{\consigna
Esta subsección debe definir el objetivo puntual del proyecto, el resultado esperable, de forma más específica que la definición previa, en cuanto a qué esperar. Debe especificarse si se espera un prototipo o un producto final, y con qué calidad. Si el resultado se compone de varias partes, deben especificarse todas.\\
El objetivo aquí descripto debe cumplir con el paradigma S.M.A.R.T. El lector debe entender completamente qué esperar como resultado del proyecto.
}

El objetivo del proyecto es construir un sistema portátil y funcional compuesto por un parche sensor descartable de biomarcadores en sudor, un módulo electrónico reutilizable de lectura y comunicación que se acopla al parche sensor, una pulsera reutilizable de señales fisiológicas y una aplicación móvil, que permita realizar un seguimiento de marcadores biométricos, identificar factores probablemente asociados a las variaciones observadas y estimar la evolución de la glucosa en dos horizontes de predicción: corto plazo (30 a 60 minutos) y horizonte extendido (1 a 8 horas). La identificación de factores se limitará a un conjunto finito y predefinido, como la alimentación, el estrés y la actividad física, sin establecer causalidad clínica. El sistema deberá presentar las mediciones, las estimaciones y los avisos correspondientes para facilitar el seguimiento cotidiano del usuario; el desempeño de sus funciones deberá verificarse mediante pruebas de aceptación.

Como resultado se entregarán cuatro componentes integrados: una aplicación móvil, un parche sensor descartable, un módulo electrónico reutilizable de lectura y comunicación, y una pulsera reutilizable. La aplicación deberá vincularse con el módulo y la pulsera, recibir y presentar sus datos e implementar las funciones de análisis, historial, avisos y ayuda comprometidas para la versión. El parche descartable comprenderá el sustrato de papel, sus zonas con reactivos y el soporte adhesivo; se establece como objetivo reemplazarlo cada 7 días de uso, sujeto a validar la estabilidad de los reactivos, la adhesión y el desempeño. El módulo electrónico reutilizable se acoplará al parche e incluirá los elementos de lectura, procesamiento, comunicación, firmware y alimentación que se definan. La pulsera reutilizable incluirá sus sensores, electrónica, firmware y alimentación. Los tres componentes físicos se entregarán protegidos en un paquete que los identifique e incluya la lista de contenido, los accesorios y consumibles definidos, y las instrucciones de puesta en marcha, conexión con la aplicación, mantenimiento de los elementos reutilizables y recambio del parche. La aplicación móvil no requiere reemplazo físico.

%\noindent\textbf{Pendiente de definir:} el plazo de desarrollo del producto, la cantidad de kits que se producirán, las plataformas compatibles con la aplicación, los umbrales de desempeño para la aceptación y si el resultado será un prototipo funcional de evaluación o un producto final comercial. Estas decisiones no se vinculan con la fecha de entrega de la actividad académica.


\subsection{Entregables mensurables}
{\consigna
Esta subsección debe definir cuáles serán los resultados esperables para verificar el éxito del proyecto.\\
Se debe especificar concretamente cómo el verificador puede constatar, al final del proyecto, que se cumplieron los objetivos del mismo. Esta descripción da certidumbre a la vez que ayuda a especificar mejor los requerimientos del proyecto.
}


Se considerará cumplido el objetivo del proyecto con la entrega y aceptación de los siguientes cuatro componentes:

\begin{enumerate}
    \item \textbf{Aplicación instalable:} deberá permitir la vinculación por BLE con el módulo electrónico reutilizable y la pulsera, la recepción y visualización de datos, la consulta del historial, las funciones de análisis comprometidas para la versión, la presentación de avisos acordes con el uso previsto y el acceso a la ayuda. La aplicación móvil no requiere reemplazo físico. Las estimaciones deberán informar su horizonte temporal y su incertidumbre. La aceptación verificará la instalación, la vinculación, la transmisión, el historial y la respuesta ante fallas; el desempeño de cada modelo se evaluará por separado conforme a los umbrales pendientes de definir.
    \item \textbf{Parche sensor descartable:} parche de papel con sus zonas con reactivos y soporte adhesivo para la medición de los biomarcadores previstos en sudor. El objetivo es reemplazar la parte no reutilizable del parche cada 7 días de uso.
    \item \textbf{Módulo electrónico reutilizable:} módulo acoplable al parche sensor, con lector óptico, front-end electroquímico, microcontrolador, BLE, firmware y alimentación, según el diseño detallado. La aceptación verificará el acoplamiento con el parche, la adquisición y procesamiento de las lecturas, la transmisión de datos a la aplicación, las tolerancias que se definan y una autonomía objetivo de al menos 7 días de uso normal.
    \item \textbf{Pulsera reutilizable:} dispositivo con PPG, sensor de temperatura, IMU, BLE, firmware y alimentación, según el diseño detallado. La aceptación verificará la adquisición y transmisión de las señales previstas, la integración con la aplicación, las tolerancias que se definan y una autonomía objetivo de al menos 7 días de uso normal. 
\end{enumerate}

Los tres componentes físicos se entregarán protegidos en un paquete identificado, con lista de contenido, accesorios y consumibles definidos, guía rápida, manual y guía de solución de problemas. Sus instrucciones deberán cubrir la puesta en marcha, la conexión con la aplicación, el mantenimiento del módulo y la pulsera, y el recambio del parche descartable. El paquete es una condición de presentación y entrega, no un quinto componente del sistema.

%****************    EVALUACION DEL PROYECTO  **********%

